\SourcePage{251}{254}
\section{The spin operator}

In the remainder of this chapter we shall not be concerned with the coordinate
dependence of wave functions. Thus, when discussing how
$\psi(x,y,z;\sigma)$ behaves under a rotation of the coordinate system, we
may suppose that the particle is at the origin. Its coordinates are then
unchanged by the rotation, and the result describes specifically the
dependence of $\psi$ on the spin variable $\sigma$.

The variable $\sigma$ differs from ordinary variables (coordinates) in being
discrete. The most general linear operator acting on functions of a discrete
variable may be written
\begin{equation}
 (\hat f\psi)(\sigma)=\sum_{\sigma'}f_{\sigma\sigma'}\psi(\sigma'),
 \tag{55.1}
\end{equation}
where $f_{\sigma\sigma'}$ are constants. The parentheses around $\hat f\psi$
emphasize that the spin argument following them belongs to the function
produced by $\hat f$, not to the original function $\psi$. The quantities
$f_{\sigma\sigma'}$ are readily seen to be the matrix elements of the
operator as ordinarily defined by (11.5)\footnote{Notice that the indices on
the matrix elements on the right of (55.1) occur in an order which is, in a
certain sense, the reverse of the usual order in (11.11).}. Coordinate
integration in (11.5) is now replaced by summation over the discrete variable,
so that
\begin{equation}
 f_{\sigma_2\sigma_1}=\sum_\sigma
 \psi_{\sigma_2}^*(\sigma)[\hat f\psi_{\sigma_1}(\sigma)].
 \tag{55.2}
\end{equation}

\SourcePage{252}{255}
Here $\psi_{\sigma_1}(\sigma)$ and $\psi_{\sigma_2}(\sigma)$ are
eigenfunctions of $\hat s_z$ belonging to $s_z=\sigma_1$ and
$s_z=\sigma_2$. Each represents a state with a definite $s_z$, so only one
wave-function component is nonzero\footnote{More precisely one should write
$\psi_{\sigma_1}(\sigma)=\psi(x,y,z)\delta_{\sigma_1\sigma}$; the coordinate
factors irrelevant here have been omitted in (55.3).

We emphasize once again the need to distinguish the specified eigenvalue of
$s_z$ ($\sigma_1$ or $\sigma_2$) from the independent variable $\sigma$.
This is precisely the reason for the difference between (11.11) and (55.1).}:
\begin{equation}
 \psi_{\sigma_1}(\sigma)=\delta_{\sigma\sigma_1},\qquad
 \psi_{\sigma_2}(\sigma)=\delta_{\sigma\sigma_2}.
 \tag{55.3}
\end{equation}
By (55.1),
\[
 (\hat f\psi_{\sigma_1})(\sigma)=
 \sum_{\sigma'}f_{\sigma\sigma'}\psi_{\sigma_1}(\sigma')=
 \sum_{\sigma'}f_{\sigma\sigma'}\delta_{\sigma'\sigma_1}
 =f_{\sigma\sigma_1}.
\]
Substitution of this result and $\psi_{\sigma_2}(\sigma)$ in (55.2) makes
that equation an identity, proving the assertion.

Operators acting on functions of $\sigma$ can therefore be represented by
matrices of order $2s+1$. In particular, this applies to the spin operator
itself, whose action is, by (55.1),
\begin{equation}
 (\hat{\mathbf s}\psi)(\sigma)=
 \sum_{\sigma'}\mathbf s_{\sigma\sigma'}\psi(\sigma').
 \tag{55.4}
\end{equation}
As stated at the end of \S~54, the matrices $\hat s_x$, $\hat s_y$, and
$\hat s_z$ are the matrices $L_x$, $L_y$, and $L_z$ obtained in \S~27, with
$L,M$ replaced by $s,\sigma$:
\begin{equation}
 \begin{aligned}
 (s_x)_{\sigma,\sigma-1}&=(s_x)_{\sigma-1,\sigma}
   =\frac12\sqrt{(s+\sigma)(s-\sigma+1)},\\
 (s_y)_{\sigma,\sigma-1}&=-(s_y)_{\sigma-1,\sigma}
   =-\frac i2\sqrt{(s+\sigma)(s-\sigma+1)},\\
 (s_z)_{\sigma\sigma}&=\sigma.
 \end{aligned}
 \tag{55.5}
\end{equation}
This defines the spin operator.

In the particularly important case $s=1/2$, $\sigma=\pm1/2$, these are
two-rowed matrices. They are written
\begin{equation}
 \hat{\mathbf s}=\frac12\hat{\boldsymbol\sigma},
 \tag{55.6}
\end{equation}

\SourcePage{253}{256}
where\footnote{When matrices are written in the form (55.7), their rows and columns are indexed by the values of $\sigma$, with the row number corresponding to the first index of a matrix element and the column number to the second. Here these indices take the values $+1/2$ and $-1/2$. According to (55.4), applying the operator means multiplying the $\sigma$th row of the matrix by the wave-function components arranged as the column
$\displaystyle\psi=\binom{\psi(1/2)}{\psi(-1/2)}$.
}
\begin{equation}
 \hat\sigma_x=\begin{pmatrix}0&1\\1&0\end{pmatrix},\qquad
 \hat\sigma_y=\begin{pmatrix}0&-i\\i&0\end{pmatrix},\qquad
 \hat\sigma_z=\begin{pmatrix}1&0\\0&-1\end{pmatrix}.
 \tag{55.7}
\end{equation}
The matrices (55.7) are called the \emph{Pauli matrices}. As it should be for
a matrix defined using eigenfunctions of $s_z$ itself,
$\hat s_z=\frac12\hat\sigma_z$ is diagonal\footnote{Using the same letter for
the spin projection and for the Pauli matrices cannot cause confusion: the
Pauli matrices carry a hat.}.

We record several characteristic properties of the Pauli matrices. Direct
multiplication of (55.7) gives
\begin{equation}
 \hat\sigma_y\hat\sigma_z=i\hat\sigma_x,\qquad
 \hat\sigma_z\hat\sigma_x=i\hat\sigma_y,\qquad
 \hat\sigma_x\hat\sigma_y=i\hat\sigma_z.
 \tag{55.8}
\end{equation}
Combining these with the general commutation rules (54.1), we find
\begin{equation}
 \hat\sigma_i\hat\sigma_k+\hat\sigma_k\hat\sigma_i=0\quad(i\ne k),
 \tag{55.9}
\end{equation}
so the Pauli matrices anticommute. These relations immediately give the useful
formulas
\begin{equation}
 \hat{\boldsymbol\sigma}^{2}=3,\qquad
 (\hat{\boldsymbol\sigma}\mathbf a)
 (\hat{\boldsymbol\sigma}\mathbf b)=
 \mathbf a\mathbf b+i\hat{\boldsymbol\sigma}(\mathbf a\times\mathbf b),
 \tag{55.10}
\end{equation}
where $\mathbf a$ and $\mathbf b$ are arbitrary vectors\footnote{Terms on the
right-hand sides of (55.8)--(55.10) that do not depend on
$\hat{\boldsymbol\sigma}$ are, of course, to be understood as constants
times the two-rowed unit matrix.}. By these relations, any scalar polynomial
in the matrices $\hat\sigma_i$ reduces to terms independent of
$\hat{\boldsymbol\sigma}$ and terms linear in it. Hence every scalar function
of $\hat{\boldsymbol\sigma}$ reduces to a linear function (see problem 1).
Finally, the traces (sums of diagonal components) of the Pauli matrices and
their products are
\begin{equation}
 \Tr\hat\sigma_i=0,\qquad \Tr(\hat\sigma_i\hat\sigma_k)=2\delta_{ik}.
 \tag{55.11}
\end{equation}

\SourcePage{254}{257}
The following sections examine spin wave functions in detail, including their
behavior under arbitrary coordinate rotations. We note here at once an
important property: their behavior under rotations about the $z$ axis.

Make an infinitesimal rotation $\delta\varphi$ about $z$. In terms of angular
momentum, here spin, its operator is $1+i\delta\varphi\hat s_z$. The functions
$\psi(\sigma)$ therefore become $\psi(\sigma)+\delta\psi(\sigma)$, with
\[
 \delta\psi(\sigma)=i\delta\varphi\hat s_z\psi(\sigma)
 =i\sigma\psi(\sigma)\delta\varphi.
\]
Writing this as $d\psi/d\varphi=i\sigma\psi(\sigma)$ and integrating, we find
that a finite rotation through $\varphi$ gives
\begin{equation}
 \psi'(\sigma)=e^{i\sigma\varphi}\psi(\sigma).
 \tag{55.12}
\end{equation}
For $\varphi=2\pi$, the multiplier is $e^{2\pi i\sigma}$, the same for every
$\sigma$ and equal to $(-1)^{2s}$, since $2\sigma$ always has the same parity
as $2s$. Thus a complete rotation of the coordinate system about $z$ restores
the wave functions of an integral-spin particle, but reverses the sign of
those of a half-integral-spin particle.

\ProblemHeading

1. Reduce an arbitrary function of the scalar
$a+\mathbf b\hat{\boldsymbol\sigma}$, which is linear in the Pauli matrices,
to a linear function.

\SolutionHeading

To find the coefficients in
\[
 f(a+\mathbf b\hat{\boldsymbol\sigma})=A+B\mathbf b\hat{\boldsymbol\sigma},
\]
choose $z$ along $\mathbf b$. The eigenvalues of
$a+\mathbf b\hat{\boldsymbol\sigma}$ are $a\pm b$, and those of the
corresponding operator $f(a+\mathbf b\hat{\boldsymbol\sigma})$ are
$f(a\pm b)$. Hence
\[
 A=\frac12[f(a+b)+f(a-b)],\qquad
 B=\frac1{2b}[f(a+b)-f(a-b)].
\]

2. Find the values of the scalar product $\mathbf s_1\mathbf s_2$ of the
spin-$1/2$ angular momenta of two particles in states for which the total spin
$\mathbf S=\mathbf s_1+\mathbf s_2$ has a definite value, 0 or 1.

\SolutionHeading

The general addition formula (31.3) for any two angular momenta gives
\[
 \mathbf s_1\mathbf s_2=1/4\quad\text{for }S=1,\qquad
 \mathbf s_1\mathbf s_2=-3/4\quad\text{for }S=0.
\]

3. Which powers of the operator $s_z$ for arbitrary spin $s$ are independent?

\SolutionHeading

The operator
\[
 (\hat s_z-s)(\hat s_z-s+1)\cdots(\hat s_z+s),
\]
formed from the differences between $\hat s_z$ and every possible eigenvalue
of $s_z$, gives zero on every wave function and is therefore itself
\SourcePageJoin{255}{258}zero. Consequently $(\hat s_z)^{2s+1}$ is
expressible through lower powers of $\hat s_z$; only the powers from 1 through
$2s$ are independent.
